\documentclass{article}
\usepackage{graphicx} % Required for inserting images
\usepackage{amsmath}
\usepackage{amsthm}
\usepackage{amssymb}
\usepackage{mathrsfs}
\newtheorem{theorem}{Theorem}
\newtheorem{lemma}{Lemma}
\newtheorem{corollary}{Corollary}
\theoremstyle{remark}\newtheorem{remark}{Remark}
\newcommand{\T}[0]{\mathcal{T}}
\newcommand{\B}[0]{\mathcal{B}}
\newcommand{\C}[0]{\mathcal{C}}
\newcommand{\R}[0]{\mathbb{R}}
\newcommand{\Q}[0]{\mathbb{Q}}
\newcommand{\Z}[0]{\mathbb{Z}}
\newcommand{\N}[0]{\mathbb{N}}
\newcommand{\e}[0]{\epsilon}

\begin{document}

(a) Show that if $\phi$ and $(\phi \rightarrow \psi)$ are tautologies, then $\psi$ is a tautology.
\begin{proof}
    Suppose that $\phi$ and $(\phi \rightarrow \psi)$ are tautologies.

    Then for any truth assignment $v$ we have $v(\phi) = T$ and $v(\phi \rightarrow \psi) = T$.

    For $v(\phi \rightarrow \psi) = T$ to hold given that the antecedent, $\phi$, is always true, it must be that the consequent, $\psi$, is always true.

    Thus $v(\psi) = T$
\end{proof}

(b) is it the case that if $\phi$ and $\psi$ are tautologies, then $(\phi \rightarrow \psi)$ is a tautology?

\begin{proof}
    Suppose that $\phi$ and $\psi$ are tautologies.

    Let $v$ be any truth assignment. Then $v(\phi) = T$ and $v(\psi) = T$.

    Thus $v(\phi \rightarrow \psi) = T$, as $v$ respects the truth table for $\rightarrow$, which gives $T$ on $(T, T)$.

    Since $v$ was arbitrary, $(\phi \rightarrow \psi)$ is a tautology, as desired.
\end{proof}

(c) It is the case that if $(\phi \rightarrow \psi)$ and $\psi$ are tautologies, then $\phi$ is a tautology?

No. The claim is universally quantified over formulas, so we refute it by
exhibiting particular formulas $\phi$ and $\psi$ for which both hypotheses
hold and the conclusion fails.

\begin{proof}[Counterexample]
    Let $\phi$ be the propositional variable $p$ and let $\psi$ be the
    formula $(p \rightarrow p)$.

    \emph{$\psi$ is a tautology.} Let $v$ be any truth assignment. The pair
    $(v(p), v(p))$ is either $(T, T)$ or $(F, F)$, and the truth table for
    $\rightarrow$ gives $T$ on both inputs, so $v((p \rightarrow p)) = T$.
    Since $v$ was arbitrary, $\psi$ is a tautology.

    \emph{$(\phi \rightarrow \psi)$ is a tautology.} Let $v$ be any truth
    assignment. By the above, $v(\psi) = T$. The truth table for
    $\rightarrow$ gives $F$ only on the input $(T, F)$; since $v(\psi) = T$,
    this input does not occur, so $v((\phi \rightarrow \psi)) = T$. Since
    $v$ was arbitrary, $(\phi \rightarrow \psi)$ is a tautology.

    \emph{$\phi$ is not a tautology.} By Theorem 2.2, any function from the
    propositional variables to $\{T, F\}$ extends to a truth assignment, so
    in particular there is a truth assignment $v$ with $v(p) = F$. Then
    $v(\phi) = F$, so $\phi$ is not a tautology.

    Thus $(\phi \rightarrow \psi)$ and $\psi$ are tautologies while $\phi$
    is not, so the answer is no.
\end{proof}

\medskip
\noindent\emph{Remark.} The second verification used nothing about $\phi$:
whenever $v(\psi) = T$, the fatal input $(T, F)$ cannot occur. So if $\psi$
is a tautology, then $(\phi \rightarrow \psi)$ is a tautology for
\emph{every} formula $\phi$ — a strengthening of part (b). The first
hypothesis of (c) is thus a consequence of the second and places no
constraint on $\phi$ whatsoever.

\begin{proof}[my version]
    Suppose $(\phi \rightarrow \psi)$ and $\psi$ are tautologies. Let $\phi$ be a propositional variable. Since $(\phi \rightarrow \psi)$ is a tautology it's range is just $\{ T \}$. We see that the preimage of the truth assignment for $\rightarrow$ contains input $\{ F, T \}$. So, it is not necessary for $\phi$ to be a tautology for $(\phi \rightarrow \psi)$ to be a tautology.
\end{proof}

\newpage
(a) show that if $\phi$ and $(\phi \rightarrow \psi)$ are tautologies then $\psi$ is a tautology.
\begin{proof}
    Suppose $\phi$ and $(\phi \rightarrow \psi)$ are tautologies.

    So let $v$ be any truth assignment.

    By definition $v(\phi) = T$ and $v((\phi \rightarrow \psi)) = T$.

    We see that $v((\phi \rightarrow \psi)) = F$ if and only if $v(\phi) = T$ and $v(\psi) = F$.

    But since $v((\phi \rightarrow \psi)) = T$ and $v(\phi) = T$, $v(\psi) \neq F$ so $v(\psi) = T$.

    Since $v$ was arbitrary, $\psi$ is a tautology.
\end{proof}

(b) is it the case that if $\phi$ and $\psi$ are tautologies, then $(\phi \rightarrow \psi)$ is a tautology?
\begin{proof}
    Suppose $\phi$ and $\psi$ are tautologies.

    So let $v$ be any truth assignment.

    By definition $v(\phi) = T$ and $v(\psi) = T$.

    We see that $v((\phi \rightarrow \psi)) = F$ if and only if $v(\phi) = T$ and $v(\psi) = F$.

    But since $v(\psi) = T$, we have $v(\psi) \neq F$, so $v((\phi \rightarrow \psi)) = T$.
    
    Since $v$ was arbitrary, $(\phi \rightarrow \psi)$ is a tautology.
\end{proof}

(c) It is the case that if $(\phi \rightarrow \psi)$ and $\psi$ are tautologies, then $\phi$ is a tautology?

\begin{proof}
    this can only be proved through a counter example :/
\end{proof}


\end{document}

